% ID: SEL-D0070-OA-E03
% ESTADO: cotejada. Segunda pasada visual y borrador compilado revisados.
% Fuente: Selectividad/sel_2009_matematicas_aplicadas/mat apli CC SS04.pdf
% SHA-256: bdea2fa323085e0db4eacff922ba9f403f58ef4e6bb4f2b569a0b857554c9ef2
% Página: 1; opción A; ejercicio original 3.

\noindent\underline{Parte I}

Sean $A$ y $B$ dos sucesos de un experimento aleatorio tales que:
\[
P(A^C)=0.2,\quad P(B)=0.25\quad\text{y}\quad P(A\cup B)=0.85.
\]
\begin{enumerate}[label=\alph*)]
\item \textbf{(1.25 puntos)} \textquestiondown{}Son los sucesos $A$ y $B$ independientes?
\item \textbf{(0.75 puntos)} Calcule $P(A^C/B^C)$.
\end{enumerate}
\par\noindent\begin{minipage}{\linewidth}
\noindent\underline{Parte II}

\textbf{(2 puntos)} Escriba todas las muestras de tamaño 2 que, mediante muestreo aleatorio simple (con reemplazamiento), se pueden extraer del conjunto $\{8,10,12\}$ y determine el valor de la varianza de las medias de esas muestras.
\end{minipage}\par
